\documentclass[11pt]{article}
\usepackage[margin=1in]{geometry}
\usepackage{amsmath,amssymb,amsthm}
\usepackage[hidelinks]{hyperref}
\newtheorem{theorem}{Theorem}
\newtheorem{lemma}{Lemma}
\newcommand{\tr}{\operatorname{tr}}
\newcommand{\HH}{\operatorname{ht}}
\begin{document}
\begin{center}
{\Large Markov Sail Duality from Christoffel matrix paths}\\[0.5em]
MathDB 366807\\ September 12, 2026
\end{center}

\paragraph{Task 1.}
Prove the Markov Sail Duality conjecture of Evans, Veselov, and Winn
\cite[Conjecture 8.3]{EVW}. We prove an explicit formula for every
interior coefficient on the sail segments involved in that conjecture.
Throughout, $2\leq a<b$ are coprime, $P=P_{a/b}(u,v,1)$ is the
homogeneous numerator specialized at $w=1$, and
\[
 \HH(i,j)=bi+aj-ab.
\]
The critical triangle consists of integer points with $i<a$, $j<b$,
and $\HH(i,j)>0$. Coefficients retain their individual identities
under $w=1$ by homogeneity.

\begin{theorem}\label{main}
At every interior integer point $(i,j)$ on a Markov sail segment
$B_rB_{r+1}$ or $A_{r+1}A_{r+2}$, the coefficient is
\[
 \mathcal M(i,j)=4\HH(i,j).
 \tag{1}
\]
Consequently both Markov Sail Duality identities hold.
\end{theorem}

We use the established generalized Cohn-matrix formula
\cite[Theorem 10.1]{EVW}, with its Christoffel-word realization
\cite[Sections 2.1.2, 2.7.2 and Theorem 2.7.5]{E}.
Put $S=1+u+v$ and
\[
 M=\begin{pmatrix}v&1\\uv&1+u\end{pmatrix},\qquad
 N=\begin{pmatrix}1+v&1\\uv&u\end{pmatrix}.
 \tag{2}
\]
If $C$ is the product along the lower Christoffel word with $b$
letters $M$ and $a$ letters $N$, then
\[
 \tr C=SP.
 \tag{3}
\]
Indeed, after $z=1$, $u=x^2$, $v=y^2$, and conjugation by
$\operatorname{diag}(x,1)$, the source's initial matrices are $M$
and $MN$. Its modified Christoffel word replaces a horizontal-vertical
pair by the latter matrix, giving precisely (3).

\begin{lemma}\label{count}
For $0<h<a$, the coefficient of the unique critical-triangle
monomial of height $h$ in $\tr C$ receives a contribution $4h$
from matrix paths with exactly one transition of each type.
Every path with $k\geq1$ transitions of each type has minimum
monomial height at least $k$.
\end{lemma}
\begin{proof}
Expand the trace as a sum over cyclic states in $\{0,1\}$, one
matrix entry per letter. A $0\to1$ entry is $1$, and a $1\to0$
entry is $uv$. The diagonal entries have minimum monomials
\[
 M_{00}=v,\quad M_{11}:1,\qquad
 N_{00}:1,\quad N_{11}=u.
 \tag{4}
\]
Choosing the other summand in a diagonal entry increases height by
$a$ or $b$.

Let $L=a+b$. At the gaps of the Christoffel word, let $H$ be
$b$ times the number of preceding $N$'s minus $a$ times the number
of preceding $M$'s. Its $L$ values are the integers
$-L+1,\ldots,0$, each once. One direct construction starts at zero
and subtracts $a$ if the result stays in this interval, otherwise
adds $b$. These steps give the lower Christoffel path; the values
are distinct because each step is $-a$ modulo $L$ and
$\gcd(a,L)=1$. A letter $M$ starts at a height in $[-b+1,0]$;
a letter $N$ starts in $[-L+1,-b]$.

Consider a single state-1 arc, beginning with a $0\to1$ transition
at letter $r$ and ending with a $1\to0$ transition at a distinct
letter $s$, in cyclic order. Let $m_r,m_s$ be $1$ for an $M$
and $0$ for an $N$. If $H_{r-1},H_{s-1}$ are the gap heights just
before these letters, its minimum monomial has height
\[
 h=(H_{s-1}-a m_s)-(H_{r-1}-b m_r)+a.
 \tag{5}
\]
To check this, if the arc's interior has $n$ letters $N$ and its
whole arc, including endpoints, has $m$ letters $M$, its monomial
is $u^{n+1}v^{b-m+1}$. Its height is $b(n+1)-a(m-1)$,
which gives (5). The total word has net height zero, so the same
calculation holds for an arc crossing its chosen initial gap.

Write $U_s=H_{s-1}-a m_s$ and $V_r=H_{r-1}-b m_r$.
Their multisets are the unions, with multiplicity, of
\[
\begin{split}
 U &: [-L+1,-a]\ \cup\ [-L+1,-b],\\
 V &: [-2b+1,-b]\ \cup\ [-L+1,-b].
\end{split}
 \tag{6}
\]
In particular, (5) is always at least $1$. For $0<h<a$, it is
equivalent to
\[
 (U_s+L-1)+(-b-V_r)=h-1.
\]
There are $h$ nonnegative pairs on the left. Every participating
value in each multiset has multiplicity two, so there are $4h$
ordered endpoint pairs. None has $r=s$: that choice would give
height $b$ for an $M$ or $a$ for an $N$.

These monomials have both exponents positive. Height $h<a\leq b$
therefore forces $i<a$ and $j<b$. For a given $h$, coprimality
gives at most one such pair of exponents, since the difference of
two solutions of $bi+aj=ab+h$ is a multiple of $(a,-b)$.
Thus the $4h$ paths contribute to the same coefficient.

Finally, a path with $k$ state-1 arcs is obtained from the all-zero
state path by making those disjoint local changes. Its minimum
height is the sum of the $k$ single-arc heights in (5), each
positive. This proves the last assertion too.
\end{proof}

\begin{lemma}\label{sail}
An interior point on the lower integer hull above the segment
joining $O=(0,b)$ and $T=(a,0)$, of height $0<h<a$, receives
no contribution from paths with two or more transitions of each type.
\end{lemma}
\begin{proof}
Such a point $Q$ minimizes a linear function $\ell(i,j)=\alpha i+
\beta j$, where $\alpha,\beta\geq0$ and not both zero, over the
integer points of $[0,a]\times[0,b]$ strictly above $OT$.
Write $d=\ell(Q)$. Then
\[
 \min\{\ell(O),\ell(T)\}<d.
 \tag{7}
\]
For example, if $\alpha a\leq\beta b$, the inequality
$i/a+j/b>1$ gives $\ell(i,j)>\alpha a$ when both coefficients
are positive. If a coefficient is zero, (7) follows from $i,j>0$.
The other case is symmetric.

Suppose a path with $k\geq2$ arcs contributes at $Q$.
There can be no nonminimum diagonal choice, since it would add at
least $a$ to the already positive height. For each individual
state-1 arc, make that change alone to the all-zero path and let
$Q_r$ be the resulting exponent point. Then
\[
 Q=O+\sum_{r=1}^k(Q_r-O),\qquad
 h=\sum_{r=1}^k\HH(Q_r).
 \tag{8}
\]
By Lemma \ref{count}, every summand height is positive and less
than $a$, and both exponents are positive. Thus each $Q_r$ is
itself a critical-triangle integer point, so $\ell(Q_r)\geq d$.
If $\ell(O)<d$, applying $\ell$ to (8) yields
$d-\ell(O)\geq k(d-\ell(O))$, a contradiction.

If instead $\ell(T)<d$, make each individual state-0 arc alone
to the all-one path. The resulting points $R_r$ likewise have
positive heights summing to $h$ and positive exponents: each is
a path with one transition of each type. They are therefore in
the critical triangle, and
\[
 Q=T+\sum_{r=1}^k(R_r-T).
\]
The same argument with $T$ gives a contradiction. Equation (7)
ensures that one of these arguments applies.
\end{proof}

\begin{proof}[Proof of Theorem \ref{main}]
The source identifies the combined Markov sail with the lower
integer hull used in Lemma \ref{sail}; see \cite[pages 20--21]{EVW}.
Write $b/a=[c_1,\ldots,c_n]$, with convergents $p_k/q_k$,
$p_{-1}=0,q_{-1}=1,p_0=1,q_0=0$. Its vertices in the
numerator's coordinate frame are
\[
 A_r=(q_{2r-1},b-p_{2r-1}),\qquad
 B_r=(a-q_{2r},p_{2r}).
 \tag{9}
\]
The positive Euclidean remainders $|bq_k-ap_k|$ strictly decrease
from $a$ at $k=0$ to zero at the final convergent. In particular,
$\HH(A_1)=b-c_1a<a$, and $\HH(B_0)=a$.
Height is affine on a segment. Consequently every interior
critical-triangle point on $B_rB_{r+1}$ or $A_{r+1}A_{r+2}$ has
height strictly between $0$ and $a$; the boundary point $B_0$
is excluded.

Lemmas \ref{count} and \ref{sail} show that its coefficient in
$\tr C$ is $4h$. Paths with no transitions have monomials
divisible by $v^b$ or $u^a$, so do not contribute.
Equation (3) gives the same coefficient in $P$: the possible
contributions from $uP$ and $vP$ would have predecessor heights
$h-b<0$ and $h-a<0$, outside the established Newton polygon
\cite[Theorem 3.2]{EVW}. This proves (1).

For each identity whose opposite vertex is interior, the primitive
direction along $B_rB_{r+1}$ is
$(-q_{2r+1},p_{2r+1})$, by the convergent recurrence. Its height
increment is $-\HH(A_{r+1})$. Similarly, the primitive direction
along $A_{r+1}A_{r+2}$ is $(q_{2r+2},-p_{2r+2})$, with height
increment $-\HH(B_{r+1})$. Applying (1) gives exactly
\[
 d(B_rB_{r+1})=-\mathcal M(A_{r+1}),\qquad
 d(A_{r+1}A_{r+2})=-\mathcal M(B_{r+1}).
\]
As in the source, $\mathcal M$ and differences here concern
interior points only. No formula is asserted for the excluded
boundary coefficients, and a segment with fewer than two interior
integer points supplies no adjacent-difference condition.
The Fibonacci cases $a=1$ have empty critical triangles.
\end{proof}

\paragraph{Task 2.}
Determine the coefficients on the initial sail segment $A_0A_1$,
including the points with height at least $a$.

\paragraph{Task 3.}
Count matrix paths with exactly two transition pairs and obtain
explicit formulas for the next layer of coefficients above the sail.

\paragraph{Task 4.}
Relate the transition-pair decomposition to local operations on
weighted snake-graph perfect matchings.

\paragraph{Task 5.}
Use the path expansion to investigate log-concavity along lines
in the full Markov Newton polygon.

\begin{thebibliography}{9}
\raggedright
\bibitem{EVW} S. J. Evans, A. P. Veselov, and B. Winn,
\emph{Arithmetic and geometry of Markov polynomials},
arXiv:2501.14882v3 (2025); current 28-page journal-hosted manuscript.
\url{https://armj.math.stonybrook.edu/Articles/250708-Winn/article.pdf}.
\bibitem{E} S. J. Evans, \emph{Quantum numbers and Markov polynomials},
doctoral thesis, Loughborough University, December 2025;
deposited February 9, 2026.
\url{https://doi.org/10.26174/thesis.lboro.31250974.v1}.
\end{thebibliography}
\end{document}
